\scanpage{53}
\mspage{9}
% archive sheet number: 52
\uline{Proposition} Soient $k$ un corps, $k'$ une extension
\margin{\rule[-4\baselineskip]{0.35pt}{4\baselineskip}}\uline{primaire} de $k$, $G$ un groupe \inserted{\unsure{comm.}} de type fini sur
$k$, $G'=G\otimes_k k'$, $H'$ un ss-groupe fermé de
$G'$, tel que $\deleted{H'_{0}}\inserted{{H'}^{\,(0)}_e}$ (la composante connexe \inserted{de $e$}) \deleted{\unsure{de la section}})
soit rigide. Alors il existe une \deleted{ss-groupe}
extension \uline{finie} \uline{radicielle} $k_1$ de $k$, \deleted{telle que},
\deleted{\unsure{purement}} et un ss-groupe \inserted{\unsure{réduct't|réduitt} fermé} $H_1$ de $G_1=G\otimes_k k_1$,
tels que $H'_1=H'\otimes_k k_1$ \deleted{\unsure{soit isom. à}} \deleted{\unsure{soit}} \deleted{\unsure{égal à}} $H_1\otimes_{k_1}k'_1$
(où $k'_1=k'\otimes_k k_1$) \deleted{\unsure{et|ait}} \inserted{aient} $=$ ensemble sous-jacent.
% slanted
\margin{? N.B. Si on veut \uline{\unsure{majorer|imposer}}, $H'$ \unsure{par un} $H$ \unsure{défini} \unsure{on|par} le \unsure{pourra|passage} au \unsure{moins} pour une extension \unsure{finie} radicielle $k_1$}
\uline{Cor. 0} Si $k'/k$ est \deleted{\unsure{sép.}} primaire séparable, \deleted{alors $H'_1\supset H_1\otimes_{k_1}k'_1$.}

% in the left margin beside Cor. 0 and Corollaire 1
\diagram{\begin{tikzcd}[ampersand replacement=\&,column sep=small,row sep=small]
k' \arrow[r,dash]\arrow[d,dash] \& k'_1 \arrow[d,dash]\\
k\arrow[r,dash] \& k_1
\end{tikzcd}}

\uline{Corollaire 1} Il existe une \deleted{\unsure{sous-schéma}}
partie fermée $H$ de $G$ telle que
$|H'|={}$ image inverse de $H$

\uline{Corollaire 2} Si $k$ est parfait, il existe un
ss-groupe $H$ de $G$ tel que $H'\supset H\otimes_k k'$ et
$H'$ a $=$ ens. sous-jacent.

\uline{Corollaire 3} Soit \inserted{$T'\longrightarrow S$ de type fini surjectif},
\deleted{$T$ de type fini sur $S$,}
à fibres relt \deleted{\unsure{conn.}} irréductibles, $S'\longrightarrow S$
% slanted
\margin{\deleted{\unsure{hyp. plus gén.} \uline{\unsure{sans vérif.}}}}
\deleted{ét. est universellement un \unsure{hom.} \unsure{?} \unsure{quotient} topologique,}
$G$ \deleted{\unsure{$S$-groupe}} \inserted{schéma en groupes \unsure{comm.}} de type fini sur $S$,
$H'$ un ss-schém. en groupes de $G'=G\times_S S'$, \inserted{sur $S'$}, \unsure{on}
suppose : \unsure{a)} $H'$ ouvert sur $S'$ \unsure{(b)} les comp.
conn. des fibres de $H'$ aux pts gén. des fibres de $S'/S$ sont
\inserted{\unsure{rigides|rigide}}
